% ============================================================
%  3. Problem Formulation
% ============================================================
\section{Problem Formulation}
\label{sec:problem}

\paragraph{Clean NER.}
Let $\bw = (w_1, \ldots, w_n)$ be a sentence of $n$ words.  NER assigns each
word a label $y_i \in \calY$ from a BIO tag set:
\[
  \calY = \bigl\{
    \texttt{O},\;
    \texttt{B-PER},\;\texttt{I-PER},\;
    \texttt{B-ORG},\;\texttt{I-ORG},\;
    \texttt{B-LOC},\;\texttt{I-LOC},\;
    \texttt{B-MISC},\;\texttt{I-MISC}
  \bigr\},
  \quad |\calY| = 9.
\]
\noindent where \texttt{O} denotes a non-entity token; \texttt{B-X} marks the \emph{beginning} of an entity of type X; \texttt{I-X} marks a \emph{continuation} token inside an entity of type X; and entity types are PER (person), ORG (organisation), LOC (location), and MISC (miscellaneous). The total tag set size is $|\calY|{=}9$.

The goal is to learn $f_\theta : \bw \mapsto (y_1,\ldots,y_n)$ that maximises
entity-level F1 on the CoNLL-2003 test set.

\paragraph{Noisy NER.}
A stochastic character-level perturbation operator
$\calN : \bw \mapsto \tilde{\bw}$ is applied to the input sentence, where
$\tilde{\bw} = (\tilde{w}_1, \ldots, \tilde{w}_m)$ with $m \neq n$ in
general.  Four perturbation types are supported:
\begin{align}
  \calN_{\text{sub}}   &: c \;\leftarrow\; \text{Uniform}(\text{visually-similar}(c))
                        \quad \text{with prob.}\ p_{\text{sub}}, \label{eq:nsub}\\
  \calN_{\text{ins}}   &: \text{insert random } c' \text{ at position } k
                        \quad \text{with prob.}\ p_{\text{ins}}, \label{eq:nins}\\
  \calN_{\text{del}}   &: \text{delete character at position } k
                        \quad \text{with prob.}\ p_{\text{del}}, \label{eq:ndel}\\
  \calN_{\text{space}} &: \text{insert space at position } k \text{ within word}
                        \quad \text{with prob.}\ p_{\text{space}}. \label{eq:nspace}
\end{align}
\noindent where $c$ is the original character at position $k$; $\text{visually-similar}(c)$ is a predefined set of characters typographically close to $c$ (e.g.\ `0' for `O', `1' for `l'); $c'$ is a randomly sampled character; and $p_{\text{sub}}, p_{\text{ins}}, p_{\text{del}}, p_{\text{space}} \in [0,1]$ are the per-character perturbation probabilities applied independently. The composite operator applies all four perturbations independently per
character in sequence.  Because $\calN_{\text{space}}$ can split one word into
two, the noisy word count $m$ may differ from $n$, requiring \emph{label
projection} (Section~\ref{subsec:label_proj}).

The objective in the noisy setting is to train $f_\theta$ such that accuracy
on $\tilde{\bw}$ degrades as little as possible relative to $\bw$.

\paragraph{Scope.}
We do not assume the noise level is known at test time; the model must be
robust across the full range of perturbations without noise-level
conditioning.
